% !TEX program = pdflatex
\documentclass[11pt]{article}

% Packages
\usepackage[margin=1in]{geometry}
\usepackage{amsmath,amssymb,amsfonts}
\usepackage{bm}
\usepackage{physics}
\usepackage{siunitx}
\usepackage{graphicx}
\usepackage{hyperref}
\usepackage{xcolor}
\hypersetup{colorlinks=true,linkcolor=blue,citecolor=blue,urlcolor=blue}

% Macros
\newcommand{\UoneEM}{\mathrm{U}(1)_{\rm EM}}
\newcommand{\hc}{\hbar c}
\newcommand{\GE}{G_E}
\newcommand{\GN}{G_E^n}
\newcommand{\GD}{G_D}
\newcommand{\Tr}{\mathrm{Tr}}
\newcommand{\Lag}{\mathcal{L}}

\title{Isospin\mbox{\textendash}Violating Scalar Mediator and a Low\mbox{\(\,Q^2\,\)} Neutron Bump\\\large Gauge\mbox{\textendash}Invariant EFT Note (v2.1)}
\author{Ryan O. Van Gelder \\ \\ Citizen Gardens}
\date{October 1, 2025}

\begin{document}
\maketitle

\begin{abstract}
We present a gauge\mbox{\textendash}invariant effective field theory (EFT) that preserves exact $\UoneEM$ while reproducing the previously proposed low\mbox{$Q^2$} ``bump'' in the neutron electric form factor $\GN$. The mechanism uses a SM\mbox{\textendash}singlet scalar $\chi$ with a gluon portal that creates a localized in\mbox{\textendash}hadron $\chi$ density and hypercharge/Yukawa portals that induce a small, isovector shift of the neutron charge distribution. The construction respects Ward identities, anomaly cancellation, and is radiatively stable under MFV alignment. The phenomenological prediction $\Delta \GN(Q^2)=\eta\,\big[e^{-\alpha_p Q^2}-e^{-\alpha_n Q^2}\big]$ and the associated experimental strategy remain unchanged.
\end{abstract}

\section{Gauge\mbox{\textendash}Invariant EFT Foundation}
\textbf{Field content:} SM $+$ real scalar $\chi$ (SM singlet). SM gauge charges and $\Lag_{\rm QCD+QED}$ unchanged.

\paragraph{Operator basis ($d\le 6$).} We consider
\begin{align}
\Lag_{\rm eff}&=\Lag_{\rm SM}+\frac{1}{2}(\partial_\mu\chi)^2-\frac{m_\chi^2}{2}\chi^2-\frac{\lambda_\chi}{4}\chi^4\\
&\quad+\frac{c_g}{\Lambda_\chi}\,\chi\,\Tr\,G_{\mu\nu}G^{\mu\nu}\quad\text{(gluon portal)}\\
&\quad+\frac{c_B}{\Lambda_\chi^2}\,\chi^2 B_{\mu\nu}B^{\mu\nu}+\frac{c_W}{\Lambda_\chi^2}\,\chi^2\,\Tr\,W_{\mu\nu}W^{\mu\nu}\quad\text{(electroweak portals)}\\
&\quad+\frac{\chi^2}{\Lambda_\chi^2}\Big(\bar q_L Y_\chi^u\,\tilde H\,u_R+\bar q_L Y_\chi^d\,H\,d_R+\text{h.c.}\Big)\quad\text{(MFV Yukawa portal)}\,.
\end{align}
\textbf{Flavor safety (MFV).} Take $Y_\chi^{u,d}=\alpha_{u,d}\,y_{u,d}$ with real $\alpha_{u,d}=\mathcal{O}(1)$ to align with SM Yukawas and suppress FCNCs. A $\chi\to-\chi$ symmetry forbids linear $\chi$--fermion couplings; the only linear term is the gluon portal (tadpole source).

\paragraph{Ward identities \\& anomalies.} EM charges remain $(+2/3,-1/3)$. All operators are gauge invariant; EM current conservation and anomaly cancellation are unmodified. Kinetic effects arise only through gauge\mbox{\textendash}invariant $\chi^2 B_{\mu\nu}B^{\mu\nu}$ and are radiatively stable.

\section{Nucleon\mbox{\textendash}Level EFT (gauge\mbox{\textendash}invariant refactor)}
We do not change SM charges. Instead, a confinement sensor $\chi$ couples to gluons so that $\langle\chi\rangle$ is enhanced in confined regions, and its effect on elastic $e$--N scattering is encoded by \emph{gauge\mbox{\textendash}invariant nucleon operators}:
\begin{align}
\Lag \supset &\sum_{N=p,n}\bar N(i\!\not\! D-M_N)N
+\sum_{N=p,n}\frac{c_D^{(N)}(\chi)}{\Lambda_\chi^2}\,\bar N\gamma^\mu N\,\partial^\nu F_{\mu\nu}
+\sum_{N=p,n}\frac{c_T^{(N)}(\chi)}{\Lambda_\chi}\,\bar N\sigma^{\mu\nu}N\,F_{\mu\nu}\,;
\end{align}
with $D_\mu=\partial_\mu+ieQ_NA_\mu$, $Q_p=+1$, $Q_n=0$. The Wilson coefficients admit a smooth expansion $c_{D,T}^{(N)}(\chi)=c_{D,T,0}^{(N)}+\hat c_{D,T}^{(N)}\,\chi+\cdots$.

\paragraph{Why the phenomenology is preserved.} The Dirac\mbox{\textendash}slope operator contributes a transverse (Ward\mbox{\textendash}identity\mbox{\textendash}compliant) vertex that shifts the \emph{electric} form\mbox{\textendash}factor slope while keeping $\GN(0)=0$. Package hadronic structure into a momentum\mbox{\textendash}dependent matching factor $f_D^{(N)}(Q^2)$ normalized to $f_D^{(N)}(0)=1$. For the neutron: $\Delta F_1^{(n)}(Q^2)=(\bar c_D^{(n)}/\Lambda_\chi^2)\,Q^2 f_D^{(n)}(Q^2)$ and $\Delta \GN=\Delta F_1^{(n)}-\tau\,\Delta F_2^{(n)}$. Taking $c_T^{(n)}$ negligible gives $\Delta \GN\approx\Delta F_1^{(n)}$ at low $Q^2$.

\paragraph{Exact reproduction of the Gaussian bump.} Choose hadronic shapes $f_D^{(p)}(Q^2)=e^{-\alpha_p Q^2}$, $f_D^{(n)}(Q^2)=e^{-\alpha_n Q^2}$. Define the isovector matching
\begin{equation}
 f_D^{(n)}(Q^2)\;\to\; \frac{e^{-\alpha_p Q^2}-e^{-\alpha_n Q^2}}{(\alpha_n-\alpha_p)Q^2},\qquad \frac{\bar c_D^{(n)}}{\Lambda_\chi^2}=\eta(\alpha_n-\alpha_p),
\end{equation}
then
\begin{equation}
\boxed{\;\Delta \GN(Q^2)=\eta\big[e^{-\alpha_p Q^2}-e^{-\alpha_n Q^2}\big]\;}
\end{equation}
with the induced radius shift $\Delta\langle r^2\rangle_n=-(3/2)\,\eta\,(a_n^2-a_p^2)$.

\paragraph{Confinement sensor realized (no charge redefinition).} The glue portal $(c_g/\Lambda_\chi)\,\chi\,G^2$ makes $\langle\chi\rangle$ track confinement (large inside hadrons, tiny asymptotically). This modulates Wilson coefficients $c_{D,T}^{(N)}(\chi)$, \emph{not} electric charges. Thus DIS (short distance) stays SM, while low\mbox{$Q^2$} elastic observables can show the predicted bump.

\section{Matching to Hadronic Operators}
In HB($\chi$)PT language, integrate out hard SM modes and retain the leading, gauge\mbox{\textendash}invariant operators that modify the electric response while preserving $\GE^N(0)$:
\begin{equation}
\Lag_{\rm had} \supset c_r^{(0,3)}\,\bar N\,(\mathbb 1,\tau^3)\,N\;\partial_\nu F^{\mu\nu}
\; +\; c_P^{(0,3)}\,\bar N\,(\mathbb 1,\tau^3)\,\sigma^{\mu\nu}N\,F_{\mu\nu}+\ldots
\end{equation}
Only the $c_r$ operators shift the \emph{electric slope} (radius) and vanish at $Q^2\to0$ due to the explicit $q^2$ factor from $\partial\!\cdot\!F$, thereby respecting the Ward identity. The $\chi$ puddle induces non\mbox{\textendash}local nucleon form factors
\begin{equation}
 c_r^{(i)} \;\to\; c_r^{(i)}\,\mathcal F_\chi^N(Q^2),\qquad \mathcal F_\chi^N(Q^2)\simeq e^{-Q^2 a_N^2/[4(\hc)^2]}.
\end{equation}
Isospin breaking arises from MFV\mbox{\textendash}aligned but non\mbox{\textendash}degenerate $\alpha_u\neq\alpha_d$ and from proton\mbox{\textendash}neutron differences in $\chi$ density, leading to distinct $a_p$ and $a_n$ and small coefficient splittings.

\section{Reproducing the Gaussian Bump}
The induced contribution to the Sachs form factor is
\begin{equation}
\Delta \GE^N(Q^2)=\kappa_N\Big[\mathcal F_\chi^N(Q^2)-\mathcal F_\chi^N(0)\Big]\simeq \kappa_N\Big[e^{-Q^2 a_N^2/[4(\hc)^2]}-1\Big],
\end{equation}
which preserves $\Delta \GE^N(0)=0$. For the neutron we parameterize the mild isovector structure by two nearby length scales, giving
\begin{equation}
\boxed{\;\Delta \GN(Q^2)=\eta\,\Big[e^{-\alpha_p Q^2}-e^{-\alpha_n Q^2}\Big],\quad \alpha_{p,n}=a_{p,n}^2/[4(\hc)^2] \;} 
\end{equation}
exactly matching the closed form used previously. The neutron mean\mbox{\textendash}square charge\mbox{\textendash}radius shift is
\begin{equation}
\Delta\langle r^2\rangle_n=-\tfrac{3}{2}\,\eta\,(a_n^2-a_p^2).
\end{equation}
\noindent\textbf{Benchmark (conservative):} $a_p=\SI{0.40}{fm}$, $a_n=\SI{0.42}{fm}$, $\eta=0.40$ $\Rightarrow$ peak few $\times 10^{-3}$ near $Q^2\sim\SI{0.1}{GeV^2}$; integral tied to $\Delta\langle r^2\rangle_n$.

% Optional figure include
% \begin{figure}[t]
%   \centering
%   \includegraphics[width=0.75\linewidth]{overlay_ge_n.png}% ensure the file is in the same dir
%   \caption{Overlay of Galster baseline, anomaly $\Delta$, and sum.}
% \end{figure}

\section{Correlated PVES Prediction (Hypercharge Portal)}
After EWSB, $B_{\mu\nu}=c_W F_{\mu\nu}-s_W Z_{\mu\nu}$. The operator $\chi^2 B_{\mu\nu}B^{\mu\nu}$ induces coherent, co\mbox{\textendash}shaped bumps in both EM and weak neutral current form factors, plus a $\gamma\leftrightarrow Z$ kinetic\mbox{\textendash}mixing term. To leading order in the small, localized correction,
\begin{equation}
\Delta \GE^{Z,n}(Q^2)\approx \rho_Z\,\Delta \GE^{\gamma,n}(Q^2),\qquad \rho_Z=\frac{s_W^2-\xi\,c_W s_W}{c_W^2},
\end{equation}
where $\xi$ (typically $|\xi|\ll1$) parameterizes the $\gamma$--$Z$ mixing contribution in the same $\chi$ background. Thus PVES at low $Q^2$ should see a same\mbox{\textendash}shape feature with a fixed normalization ratio $\rho_Z$ that is calculable from the portal coefficients.

\section{Theoretical Consistency \\& Radiative Stability}
\begin{itemize}
  \item \textbf{Ward identities:} modifications scale as $q^2$ and vanish as $Q^2\to0$; charges unchanged.
  \item \textbf{MFV alignment:} $Y_\chi^{u,d}\propto y_{u,d}$ avoids FCNCs; CP phases can be set to zero at this stage.
  \item \textbf{Power counting:} with $m_\chi\sim\mathcal{O}(1\,\mathrm{GeV})$ and $\Lambda_\chi\gtrsim$ few GeV, the bump amplitude $\eta\sim10^{-3}$--$10^{-2}$ is natural given localized $\langle\chi^2\rangle$ at hadronic scales.
\end{itemize}

\section{Experimental Strategy (unchanged)}
\begin{itemize}
  \item \textbf{Target observable:} $\Delta \GN(Q^2)$ for $Q^2\in[0.05,0.30]~\mathrm{GeV}^2$.
  \item \textbf{Binning:} 5--7 bins with $\sigma_{\rm stat}\lesssim 3\times10^{-3}$ per bin; common\mbox{\textendash}mode systematics at the $10^{-3}$ level.
  \item \textbf{Optimization:} fit $(\eta,a_p,a_n)$ with the Gaussian\mbox{\textendash}difference template.
  \item \textbf{PVES cross\mbox{\textendash}check:} a synchronized low\mbox{$Q^2$} PVES point (or short scan) to test the $\rho_Z$ normalization prediction.
\end{itemize}

\section{Quick Constraint Posture}
\begin{itemize}
  \item \textbf{LHC / quarkonium:} $\chi$ pair production via gluons with loop\mbox{\textendash}suppressed $\gamma\gamma$ decays (from $c_B$) is weakly constrained for $\Lambda_\chi$ above a few GeV and $m_\chi\gtrsim$ GeV; the amplitude required here is small.
  \item \textbf{Nuclear/atomic:} The localized, in\mbox{\textendash}hadron nature and the $q^2$ suppression protect atomic parity violation and Thomson\mbox{\textendash}limit observables.
\end{itemize}

\paragraph{Appendix A: Operator Basis \\& Normalizations}
\begin{itemize}
  \item Charge\mbox{\textendash}radius operators: $\mathcal{O}_r^{(0,3)}=\bar N\,(\mathbb 1,\tau^3)\,N\,\partial_\nu F^{\mu\nu}$. Matrix element contributes $\propto Q^2$ to $\GE$, preserving $\GE(0)$.
  \item Pauli operators: $\mathcal{O}_P^{(0,3)}=\bar N\,(\mathbb 1,\tau^3)\,\sigma^{\mu\nu}N\,F_{\mu\nu}$ (affects $G_M$ primarily; subleading leakage into $\GE$ via kinematics).
  \item EW matching: From $\chi^2 B^2$, after EWSB: $\chi^2(c_W^2 F^2-2c_W s_W F\!\cdot\!Z+s_W^2 Z^2)$. The induced non\mbox{\textendash}local nucleon response inherits the same $Q^2$ shape through the common $\chi$ density.
\end{itemize}

\paragraph{Appendix B: Benchmarks}
\begin{itemize}
  \item \textbf{Baseline:} $a_p=\SI{0.40}{fm}$, $a_n=\SI{0.42}{fm}$, $\eta=0.40$. Predicts peak $|\Delta \GN|\sim$ few $\times 10^{-3}$ around $Q^2\sim\SI{0.1}{GeV^2}$; integral tied to $\Delta\langle r^2\rangle_n=-(3/2)\,\eta(a_n^2-a_p^2)$.
\end{itemize}

\nocite{*}
\bibliographystyle{unsrt}
\bibliography{references}

\end{document}
